\documentclass{article}
\usepackage{multirow}
\usepackage{tabularx}

\usepackage{graphicx} % Required for inserting images

\usepackage[a4paper, left=1.5cm, right=1.5cm, top=1cm, bottom=2cm]{geometry}

\usepackage{float}

\newcounter{commentCount}
\newcounter{filePrg}
\newcounter{inputPrg}

\usepackage[dvipsnames]{xcolor}
\usepackage{minted}

\usepackage[many]{tcolorbox}
\tcbuselibrary{listings}
\tcbuselibrary{minted}

\usepackage{ifthen}
\usepackage{fontawesome}

\usepackage{tabularx}
\newcolumntype{\CeX}{>{\centering\let\newline\\\arraybackslash}X}%
\newcommand{\TwoSymbolsAndText}[3]{%
  \begin{tabularx}{\textwidth}{c\CeX c}%
    #1 & #2 & #3
  \end{tabularx}%
}

\newtcblisting[use counter=inputPrg, number format=\arabic]{codeInput}[3]{
  listing engine=minted,
  minted language=#1,
  minted options={autogobble,breaklines, fontsize=\footnotesize},
  listing only,
  size=title,
  arc=1.5mm,
  breakable,
  enhanced jigsaw,
  colframe=myblue,
  coltitle=White,
  boxrule=0.5mm,
  colback=white,
  coltext=Black,
  title=\TwoSymbolsAndText{\faCode}{%
  \textbf{JSON}\ifthenelse{\equal{#2}{}}{}{\textbf{: }#2}%
  }{\faCode},
  label=inputPrg:#3
}


\usepackage{boldline}
\usepackage{tikz,tcolorbox}
\usepackage{amsmath}
\usepackage[table,xcdraw]{xcolor}
\usepackage{listings}
\usepackage{array,multirow} % For customizing tables
\usepackage{booktabs} % For better horizontal lines
\usepackage{makecell}
\setlength{\parindent}{0pt}
\usepackage{siunitx}
\usepackage{tkz-tab}
\usepackage{amssymb}
\usepackage{amsmath}
\usepackage{enumitem}

\usepackage{caption}
\usepackage{float}


\newcommand{\exer}[1]{
  \section*{Exercice #1}
  \vspace{-0.5cm}
  \noindent\rule{\textwidth}{0.5pt}%
}

\newcommand{\tit}[1]{
\begin{center}
    \Large{\textbf{{#1}}}
\end{center}
}

\definecolor{commentgray}{HTML}{676160}
\definecolor{messagegreen}{HTML}{17B867}
\definecolor{myblue}{HTML}{10C2C4}

\tcbuselibrary{skins, breakable, theorems}


\newtcolorbox{prettyBox}[2]{
  enhanced,
  colback=white!90!#2,   % Background color based on the second parameter (color)
  colframe=#2!60!black,  % Frame color based on the second parameter (color)
  coltitle=white,        % Title color (white)
  fonttitle=\bfseries\Large,
  title=#1,              % Title from the first parameter
  boxrule=1mm,
  arc=0.5mm,
  drop shadow=#2!35!gray, % Drop shadow color based on the second parameter (color)
}

\usepackage{hyperref}
\hypersetup{
    colorlinks=true,
    linkcolor=blue,
    filecolor=magenta,      
    urlcolor=cyan,
    pdftitle={Overleaf Example},
    pdfpagemode=FullScreen,
    }



\lstdefinestyle{cmd}{
 basicstyle=\ttfamily,
 backgroundcolor=\color{lightgray!20},
 frame=single
}
\usepackage{subcaption}
\usepackage{minted}
\usepackage{multicol}


\begin{document}

\tit{DW 1}

\exer{1}

\begin{enumerate}
    \item UML :
\begin{figure}[H] 
  \centering 
  \includegraphics[width=0.8\textwidth]{ex1.png} 
\end{figure}

\vspace{0.25cm}

\item Modelisation oriente cle valeur avec isbn comme cle :

\begin{multicols}{2}

\vspace*{\fill}
"isbn" : "2154889522"   $\xrightarrow{\hspace{4.5cm}}$
\vspace*{\fill}

\columnbreak

\begin{codeInput}{json}{Valeur}{code01}
{
    "titre" : "Asterix Et Cleopatre",
    "auteur" : { 
        "id_auteur" : 154,
        "nom" : "Gosciny",
        "prenom" : "Rene"
    },
    "ventes" : [{
        "id_vente" : 10,
        "date" : "02/06/2017"
    }]
}
\end{codeInput}

\end{multicols}



\begin{multicols}{2}

\vspace*{\fill}
"isbn" : "2154889589"   $\xrightarrow{\hspace{4.5cm}}$
\vspace*{\fill}

\columnbreak

\begin{codeInput}{json}{Valeur}{code01}
{
    "titre" : "NoSQL",
    "auteur" : { 
        "id_auteur" : 987,
        "nom" : "Bruchez",
        "prenom" : "Rudi"
    },
    "ventes" : [
    {
        "id_vente" : 12,
        "date" : "02/06/2017"
    },
     {
        "id_vente" : 20,
        "date" : "12/09/2017"
     }
    ]
}
\end{codeInput}

\end{multicols}


\vspace{0.25cm}


\item Modelisation oriente colonne pour acceder aux livres d'un auteur :



\begin{center}
\begin{tabular}{ |c|c| } 
\hline
\multirow{4}{*}{\textbf{Auteur}} & \texttt{id\_auteur} \\
                                 & \texttt{nom}        \\
                                 & \texttt{prenom}     \\
                                    \cline{2-2}
                                 & \begin{tabular}{c|c}
                                     \multirow{3}{*}{\textbf{Livre}} & \texttt{isbn}\\
                                                                     & \texttt{titre}\\
                                                                     & \begin{tabular}{c|c}
                                                                         \hline
                                                                         \multirow{2}{*}{\textbf{Ventes}} & \texttt{id\_vente}\\
                                                                                                          & \texttt{date}\\
                                                                       \end{tabular}\\
                                 \end{tabular}\\\hline
\end{tabular}
\end{center}

\newpage

\item Modelisation oriente document :

\begin{codeInput}{json}{Centre Vente}{code01}
{  
    "id_vente" : 10,
    "date" : "02/06/2017",
    "livre" : {
        "isbn" : "2154889522",
        "titre" : "Asterix Et Cleopatre",
        "auteur" : { 
            "id_auteur" : 154,
            "nom" : "Gosciny",
            "prenom" : "Rene"
        }
    }

},
{
    "id_vente" : 12,
    "date" : "02/06/2017",
    "livre" : {
        "isbn" : "2154889589",
        "titre" : "NoSQL",
        "auteur" : { 
            "id_auteur" : 987,
            "nom" : "Bruchez",
            "prenom" : "Rudi"
        } 
    }
},
{
    "id_vente" : 20,
    "date" : "12/09/2017",
    "livre" : {
        "isbn" : "2154889589",
        "titre" : "NoSQL",
        "auteur" : { 
            "id_auteur" : 987,
            "nom" : "Bruchez",
            "prenom" : "Rudi"
        } 
    }
}
\end{codeInput}



\newpage


\begin{codeInput}{json}{Centre Livre}{code01}
{   
    "isbn" : "2154889589",
    "titre" : "NoSQL",
    "auteur" : { 
        "id_auteur" : 987,
        "nom" : "Bruchez",
        "prenom" : "Rudi"
    },
    "ventes" : [{
        "id_vente" : 12,
        "date" : "02/06/2017"
    },
     {
        "id_vente" : 20,
        "date" : "12/09/2017"
    }]
},
{   
    "isbn" : "2154889522",
    "titre" : "Asterix Et Cleopatre",
    "auteur" : { 
        "id_auteur" : 154,
        "nom" : "Gosciny",
        "prenom" : "Rene"
    },
    "ventes" : [{
        "id_vente" : 10,
        "date" : "02/06/2017"
    }]
}
\end{codeInput}

\vspace{0.5cm}


\begin{codeInput}{json}{Centre Auteur}{code01}
{
    "id_auteur" : 154,
    "nom" : "Gosciny",
    "prenom" : "Rone",
    "livres" : [{
        "isbn" : "2154889522",
        "titre" : "Asterix Et Cleopatre",
        "ventes" : [
        {
            "id_vente" : 10,
            "date" : "02/06/2017"   
        }]
    }]
},
{
    "id_auteur" : 987,
    "nom" : "Bruchez",
    "prenom" : "Rudi",
    "livres" : [{   
        "isbn" : "2154889589",
        "titre" : "NoSQL",
        "ventes" : [{
            "id_vente" : 12,
            "date" : "02/06/2017"
        },
        {
            "id_vente" : 20,
            "date" : "12/09/2017"
        }]
    }]
}
\end{codeInput}

\vspace{0.5cm}

\item La modelisation influence la complexite de la requete sur les informations qui sont profondement imbrique
    dans les documents.


\vspace{0.5cm}

\item On imbriquant un vecteur de documents d'auteurs.

\newpage
\end{enumerate}


\exer{2}
\begin{enumerate}
    \item UML :
\begin{figure}[H] 
  \centering 
  \includegraphics[width=0.8\textwidth]{ex2.png} 
\end{figure}

\vspace{0.5cm}

\item Modelisation oriente document centre sur les ue :


\begin{codeInput}{json}{Centre UE}{code01}
{
    "id" : "ue:11",
    "titre" : "Java",
    "Etudiants" : [{
        "_id" : 978,
        "nom" : "ADIMI Meriem",
        "note" : 12
    }]
},
{
    "id" : "ue:27",
    "titre" : "Base De Donnees",
    "Etudiants" : [{
        "_id" : 978,
        "nom" : "ADIMI Meriem",
        "note" : 17
    },
    {
        "_id" : 476,
        "nom" : "BELABDI Ahmed",
        "note" : 10
    }]
},
{
    "id" : "ue:37",
    "titre" : "Reseaux",
    "Etudiants" : [{
        "_id" : 978,
        "nom" : "ADIMI Meriem",
        "note" : 14
    }]
},
{
    "id" : "ue:13",
    "titre" : "Methodologie",
    "Etudiants" : [{
        "_id" : 476,
        "nom" : "BELABDI Ahmed",
        "note" : 17
    }]
},
{
    "id" : "ue:76",
    "titre" : "Conduite Projet",
    "Etudiants" : [{
        "_id" : 476,
        "nom" : "BELABDI Ahmed",
        "note" : 11
    }]
},
\end{codeInput}

\end{enumerate}

\newpage

\exer{3}

\begin{enumerate}
    \item UML :

\begin{figure}[H] 
  \centering 
  
\includegraphics[width=0.8\textwidth]{ex3.png} 
\end{figure}

\vspace{0.5cm}

\item Modelisation oriente document centre etudiant , cour :


\begin{codeInput}{json}{Centre Etudiant}{code01}
{
    "nom" : "Rabah",
    "prenom" : "Chabane Chaouche",
    "Adresse" : { 
        "numRue" : 23,
        "nomRue" : "nomrue213",
        "ville" : "Alger",
        "code postale" : 16000
    },
    "Cours" : [{
        "titre" : "Java",
        "code" : 2004,
        "credit" : 4,
        "description" : "Programmation Oriente Object",
        "pre" : [{ 
            "titre" : "Algo 1",
            "code" : 2022,
            "credit" : 6,
            "description" : "Algo & Structure Donne"
        },
        {
            "titre" : "Algo 2",
            "code" : 2023,
            "credit" : 6,
            "description" : "Structure Donne Dynamique"
        }]
    }]
}
\end{codeInput}

\newpage

\begin{codeInput}{json}{Centre Cour}{code01}
{
    "titre" : "Java",
    "code" : 2004,
    "credit" : 4,
    "description" : "Programmation Oriente Object",
    "pre" : [{ 
        "titre" : "Algo 1",
        "code" : 2022,
        "credit" : 6,
        "description" : "Algo & Structure Donne"
    },
    {
        "titre" : "Algo 2",
        "code" : 2023,
        "credit" : 6,
        "description" : "Structure Donne Dynamique"
    }],
    "etudiants" : [{
        "nom" : "Rabah",
        "prenom" : "Chabane Chaouche",
        "Adresse" : { 
            "numRue" : 23,
            "nomRue" : "nomrue213",
            "ville" : "Alger",
            "code postale" : 16000
    }
    }]
}
\end{codeInput}


\vspace{0.25cm}

\begin{prettyBox}{Modélisation}{myblue}
    \begin{itemize}
        \item \textbf{Centre Étudiant} : redondance des informations sur les cours et les prérequis. 
        Les prérequis sont fortement imbriqués. 

        \item \textbf{Centre Cours} : les prérequis ne sont pas fortement imbriques. Cependant, il y a une redondance des informations 
        concernant les étudiants inscrits.
    \end{itemize}
\end{prettyBox}\vspace{0.5cm}


\item Solution :
\begin{multicols}{2}

\begin{codeInput}{json}{Etudiant}{code01}
{
    "nom" : "Rabah",
    "prenom" : "Chabane Chaouche",
    "Adresse" : { 
        "numRue" : 23,
        "nomRue" : "nomrue213",
        "ville" : "Alger",
        "code postale" : 16000
    },
    "Cours" : [{
        "titre" : "Java",
        "code" : 2004
    }]
}
\end{codeInput}

\columnbreak

\begin{codeInput}{json}{Detail Cour}{code01}
{
    "titre" : "Java",
    "code" : 2004,
    "credit" : 4,
    "description" : "Programmation Oriente Object",
    "pre" : [{ 
        "titre" : "Algo 1",
        "code" : 2022,
        "credit" : 6,
        "description" : "Algo & Structure Donne"
    },
    {
        "titre" : "Algo 2",
        "code" : 2023,
        "credit" : 6,
        "description" : "Structure Donne Dynamique"

    }]
}
\end{codeInput}

\end{multicols}



\end{enumerate}


\end{document}
